\section{Conclusion}

\paragraph{Discussion.} This work addresses the key limitation of global $\epsilon$-robustness: its stringent requirements on classifiers. As clearly visible in Theorem~\ref{thm:bounds-total}, the parameter $\mu$ increases the upper bound on $\epsilon$ in totally ordered semimetric spaces such as floating-point numbers (Section~\ref{subsec:existence-order}). $(\mu,\epsilon)$-robustness is proven equivalent in Theorem~\ref{thm:separability} to a coverage condition in parallel to $\epsilon$-robustness being equivalent to separability. This leads to necessary and sufficient conditions for the existence of a non-constant $(\mu, \epsilon)$-classifier in the general case (Theorem~\ref{thm:existence}) and for totally ordered semimetric spaces (Corollary~\ref{cor:existence-total}). Precision-dependent upper bounds (Corollaries~\ref{cor:bounds-total} and~\ref{cor:bounds-posets}) create a connection between the underlying arithmetic and the existence of non-constant classifiers. Setting $\mu = 0$, one obtains well-known results on $\epsilon$-robustness that were, to the best of our knowledge, never collected and formally proven in a single place.

More generally, the results show that robustness is not solely a property of a classifier but also of the structure of its input domain relative to a chosen distance. By relating robustness to covering and the spacing of representable values, our study of $(\mu,\epsilon)$-robustness produced exact statements that demonstrate how the choice of domain and distance constrains the existence of robust classifiers. This perspective is relevant for neural networks, which operate on finite-precision arithmetic rather than continuous domains, a mismatch between which complicates proofs of soundness for verifiers~\cite{jia2021exploiting,cordeiro2025neural}.

\paragraph{Future work.} The next logical step is to design a verification algorithm for $(\mu,\epsilon)$-robustness. This is a significant open research problem, because the $(\mu,\epsilon)$-robustness property is in the class of $\forall\exists$-hyperproperties. Although these are studied in general context~\cite{correnson2024finding,correnson2025symbolic,finkbeiner2026concurrent}, neural network verification currently focuses on the domain of hyperproperties with $\forall$ quantifiers only~\cite{athavale2024verifying,aleksandrov2026rocq,boetius2023verifying}.
